\chapter{Instrumental Variables and Local Average Treatment Effects}
\label{chap:late}

\begin{description}
  \item[Motivation.] What can a randomized offer identify when people do not
  fully comply with their assignment?
  \item[LATE.] Define the causal effect for compliers and state the assumptions
  that identify it.
  \item[IV estimation.] Derive the Wald estimand and its moment-equation
  representation.
  \item[Applications.] Discuss where credible instruments come from and why
  exclusion and relevance require substantive arguments.
\end{description}

\section{Noncompliance and intent to treat}

Suppose access to an HPV vaccine is allocated by lottery.  Let $Z_i=1$ mean
that individual $i$ wins a free offer and let $D_i=1$ mean that she actually
receives the vaccine.  Even if $Z_i$ is randomized, compliance need not be
perfect:
\begin{itemize}
  \item some winners may decline the vaccine; and
  \item some nonwinners may obtain it elsewhere.
\end{itemize}
The difference
\[
  \E[Y_i\mid Z_i=1]-\E[Y_i\mid Z_i=0]
\]
is the \emph{intent-to-treat} (ITT) effect of the offer on the outcome.  It is
not generally the average effect of receiving the vaccine.

To describe noncompliance, define the potential treatment $D_i(z)$: whether
unit $i$ would receive treatment if the instrument were set to $z$.  Only one
of $D_i(0)$ and $D_i(1)$ is observed, because
\begin{equation}
  D_i=Z_iD_i(1)+(1-Z_i)D_i(0).
  \label{eq:late-observed-treatment}
\end{equation}
The pair $(D_i(0),D_i(1))$ defines four principal strata.

\begin{table}[tbp]
  \centering
  \caption{Compliance types for a binary instrument and treatment}
  \label{tab:late-types}
  \begin{tabular}{lcc}
    \toprule
    Type & $D_i(0)$ & $D_i(1)$\\
    \midrule
    Never-taker  & 0 & 0\\
    Complier     & 0 & 1\\
    Defier       & 1 & 0\\
    Always-taker & 1 & 1\\
    \bottomrule
  \end{tabular}
\end{table}

We never observe an individual's type directly.  Calling a unit a complier is
also application-specific: it depends on which direction of the instrument is
defined as encouragement.

\begin{definition}[Local average treatment effect]
For binary $Z_i$ and $D_i$, the local average treatment effect is
\begin{equation}
  LATE=\E[Y_i(1)-Y_i(0)\mid D_i(1)>D_i(0)].
  \label{eq:late-definition}
\end{equation}
It is the average treatment effect for compliers.
\end{definition}

If everyone is a complier, LATE equals the population ATE.  Otherwise it is a
local parameter for the subpopulation whose treatment status is changed by the
instrument.  Different valid instruments may therefore identify different
LATEs.

\section{Identification of LATE}
\label{sec:late-identification}

It helps to separate four assumptions that are sometimes bundled together.

\begin{enumerate}
  \item \textbf{Random assignment (instrument exogeneity):}
  \begin{equation}
    Z_i\indep\bigl(Y_i(0),Y_i(1),D_i(0),D_i(1)\bigr).
    \label{eq:late-random-assignment}
  \end{equation}
  \item \textbf{Exclusion:} the instrument affects the outcome only through
  treatment.  If outcomes were initially written $Y_i(d,z)$, exclusion says
  $Y_i(d,z)=Y_i(d)$ for every $d,z$.
  \item \textbf{Relevance:}
  \begin{equation}
    \E[D_i\mid Z_i=1]\neq \E[D_i\mid Z_i=0].
    \label{eq:late-relevance}
  \end{equation}
  \item \textbf{Monotonicity:}
  \begin{equation}
    D_i(1)\geq D_i(0)\quad\text{almost surely}.
    \label{eq:late-monotonicity}
  \end{equation}
  After orienting $Z_i$ as an encouragement, there are no defiers.
\end{enumerate}

Random assignment and exclusion are distinct.  A lottery can be randomized but
still violate exclusion if winning changes behavior through publicity, income,
or expectations even when treatment receipt is held fixed.

\begin{theorem}[Imbens--Angrist LATE theorem]
Under random assignment, exclusion, relevance, and monotonicity,
\begin{equation}
  LATE=
  \frac{\E[Y_i\mid Z_i=1]-\E[Y_i\mid Z_i=0]}
       {\E[D_i\mid Z_i=1]-\E[D_i\mid Z_i=0]}.
  \label{eq:late-wald-population}
\end{equation}
The numerator is the outcome ITT and the denominator is the first-stage ITT.
\end{theorem}

\begin{proof}
Consistency and exclusion give
$Y_i=D_iY_i(1)+(1-D_i)Y_i(0)$.  Using
Equation~\eqref{eq:late-observed-treatment} and random assignment,
\begin{align*}
  &\E[Y_i\mid Z_i=1]-\E[Y_i\mid Z_i=0]\\
  &\quad=\E\!\left[(D_i(1)-D_i(0))
                     (Y_i(1)-Y_i(0))\right].
\end{align*}
Under monotonicity, $D_i(1)-D_i(0)$ equals one for compliers and zero for
always- and never-takers.  Hence the last display equals
\[
  \Pr\{D_i(1)>D_i(0)\}
  \E[Y_i(1)-Y_i(0)\mid D_i(1)>D_i(0)].
\]
Similarly,
\begin{align*}
  \E[D_i\mid Z_i=1]-\E[D_i\mid Z_i=0]
  &=\E[D_i(1)-D_i(0)]\\
  &=\Pr\{D_i(1)>D_i(0)\}.
\end{align*}
Relevance makes this probability nonzero, so division yields
Equation~\eqref{eq:late-wald-population}.
\end{proof}

Monotonicity is substantive, not a mathematical convenience.  It can be
credible for a one-sided encouragement but implausible if the instrument pushes
different people in opposite directions.  Without monotonicity, the Wald ratio
mixes complier and defier effects with opposite signs.

\subsection{A linear moment representation}
\label{sec:late-linear}

For binary $Z_i$, define
\begin{align*}
  \gamma&=\E[D_i\mid Z_i=1]-\E[D_i\mid Z_i=0],\\
  \delta&=\E[Y_i\mid Z_i=1]-\E[Y_i\mid Z_i=0].
\end{align*}
The linear projections on $(1,Z_i)$ can be written
\begin{align*}
  D_i&=\gamma_0+\gamma Z_i+\nu_i,
       &\E[\nu_i]=\E[Z_i\nu_i]=0,\\
  Y_i&=\delta_0+\delta Z_i+\mu_i,
       &\E[\mu_i]=\E[Z_i\mu_i]=0.
\end{align*}
Because relevance implies $\gamma\neq0$, substituting the first equation into
the second produces
\begin{equation}
  Y_i=\beta_0+\beta_1D_i+\varepsilon_i,
  \qquad
  \E[\varepsilon_i]=\E[Z_i\varepsilon_i]=0,
  \label{eq:late-iv-linear}
\end{equation}
where $\beta_1=\delta/\gamma=LATE$.  In general
$\Cov(D_i,\varepsilon_i)\neq0$, which is precisely why OLS is not appropriate.

Equation~\eqref{eq:late-iv-linear} is a useful moment representation.  It does
not require the individual treatment effect to be constant or the conditional
mean of $Y_i$ to be globally linear in $D_i$.

\subsection{Conditional LATE}

Let $W_i$ contain pre-instrument covariates.  A conditional design assumes, for
example,
\[
  Z_i\indep\bigl(Y_i(0),Y_i(1),D_i(0),D_i(1)\bigr)\mid W_i,
\]
together with conditional exclusion, relevance, and monotonicity.  One possible
structural shorthand is $Y_i=g(D_i,W_i,U_i)$ and
$D_i=h(Z_i,W_i,V_i)$ with $Z_i\indep(U_i,V_i)\mid W_i$; importantly, the
treatment equation depends on $Z_i$, not on $D_i$ itself.

At covariate values with a nonzero conditional first stage,
\begin{equation}
  CLATE(w)=
  \frac{\E[Y_i\mid Z_i=1,W_i=w]-\E[Y_i\mid Z_i=0,W_i=w]}
       {\E[D_i\mid Z_i=1,W_i=w]-\E[D_i\mid Z_i=0,W_i=w]}.
  \label{eq:clate}
\end{equation}
Unlike the scalar unconditional LATE, $CLATE(w)$ is a function.  Flexible
methods such as \texttt{grf::instrumental\_forest()} can estimate heterogeneous
conditional effects under suitable conditions.

\section{Instrumental-variable estimation}
\label{sec:late-estimation}

The two moment equations associated with
Equation~\eqref{eq:late-iv-linear} are
\begin{align}
  \E[Y_i-\beta_0-\beta_1D_i]&=0,\label{eq:iv-moment-one}\\
  \E[Z_i(Y_i-\beta_0-\beta_1D_i)]&=0.\label{eq:iv-moment-two}
\end{align}
Eliminating the intercept gives
\begin{equation}
  \beta_1=\frac{\Cov(Y_i,Z_i)}{\Cov(D_i,Z_i)}.
  \label{eq:iv-covariance-ratio}
\end{equation}
For $p=\Pr(Z_i=1)$,
\begin{align*}
  \Cov(D_i,Z_i)
  &=p(1-p)\{\E[D_i\mid Z_i=1]-\E[D_i\mid Z_i=0]\},
\end{align*}
so the denominator is nonzero exactly when $0<p<1$ and the instrument is
relevant.  Replacing moments by sample averages yields
\begin{equation}
  \widehat\beta_{IV}=
  \frac{\sum_{i=1}^n(Y_i-\bar Y)(Z_i-\bar Z)}
       {\sum_{i=1}^n(D_i-\bar D)(Z_i-\bar Z)}.
  \label{eq:iv-sample-wald}
\end{equation}
With one binary instrument and one endogenous binary treatment, this is the
sample analog of the Wald ratio in
Equation~\eqref{eq:late-wald-population}.

\subsection{Weak instruments}

When the first-stage difference is close to zero, the ratio in
Equation~\eqref{eq:iv-sample-wald} is unstable.  Weak instruments can produce
substantial finite-sample bias and a poor normal approximation.  The familiar
``first-stage $F>10$'' rule is only a rough historical diagnostic, not a general
guarantee.  It depends on the number of instruments, the estimand, and the
desired tolerance for bias or size distortion.  Empirical work should report
the first stage and use weak-instrument-robust diagnostics or confidence sets
when weakness is a concern.

\section{Where do instruments come from?}
\label{sec:late-applications}

Instrumental variables transformed empirical economics, but a software command
cannot make an instrument credible.  Every application must defend both the
exclusion restriction and relevance.

\subsection{Randomized offers}

RCT assignment is an especially transparent instrument when actual treatment
has noncompliance.  For example, Kline and Walters (2016) use randomized offers
in the Head Start Impact Study to learn about early-childhood programs in the
presence of close substitutes.  The lottery identifies the effect for families
whose program participation changes because of the offer.

\subsection{Natural and quasi-experiments}

Nature or policy can generate plausibly exogenous variation.  Miguel, Satyanath,
and Sergenti (2004), for example, use rainfall shocks as an instrument for
economic growth when studying civil conflict in agriculturally dependent
African countries.  The exclusion argument is demanding: Sarsons (2015) shows
why rainfall may be related to conflict through channels other than local
income.  This exchange illustrates an important habit---actively search for
alternative pathways from $Z$ to $Y$.

Policy expansions in eligibility for maternal care provide another possible
instrument for service use.  Staggered implementation may create relevance,
but selective migration, concurrent programs, or direct income effects can
threaten exclusion.

\subsection{Variables excluded from the outcome equation}

When no experiment is available, researchers sometimes use a variable believed
to shift treatment but not the outcome directly.  Classic returns-to-schooling
examples include quarter of birth (Angrist and Krueger, 1991) and proximity to
college (Card, 1995).  Both require contestable institutional arguments:
season of birth and place of residence may be related to family background,
and their first-stage effects may be weak.

Other common constructions include reference-group variables and
``Hausman-type'' instruments, such as prices of the same product in other
markets.  Aggregation or geographic distance does not automatically guarantee
exogeneity; correlated demand shocks can invalidate the instrument.

\section{R implementation}
\label{sec:late-r}

The \texttt{ivreg} package uses the formula
\texttt{outcome \textasciitilde{} treatment \textbar{} instrument}.  The
following simulation contains never-takers, compliers, and always-takers but no
defiers.  The treatment effect is heterogeneous, so OLS and IV need not target
the same parameter.

\begin{rcode}
library(ivreg)

set.seed(5280)
n <- 4000
Z <- rbinom(n, 1, 0.5)
type <- sample(c("never", "complier", "always"), n,
               replace = TRUE, prob = c(0.25, 0.55, 0.20))

D0 <- as.integer(type == "always")
D1 <- as.integer(type %in% c("complier", "always"))
D <- ifelse(Z == 1, D1, D0)

U <- rnorm(n)
tau <- 1 + 0.7 * (type == "complier") + 0.4 * U
Y0 <- 0.5 + U + rnorm(n)
Y <- Y0 + D * tau

outcome_itt <- mean(Y[Z == 1]) - mean(Y[Z == 0])
first_stage <- mean(D[Z == 1]) - mean(D[Z == 0])
wald <- outcome_itt / first_stage
true_late <- mean(tau[type == "complier"])

c(outcome_itt = outcome_itt,
  first_stage = first_stage,
  wald = wald,
  true_late = true_late)

fit_iv <- ivreg(Y ~ D | Z)
summary(fit_iv, diagnostics = TRUE)
summary(lm(Y ~ D))
\end{rcode}

The companion Quarto Live page runs this example entirely in the browser.  In
real applications, diagnostics supplement rather than replace the design-based
arguments for random assignment, exclusion, monotonicity, and relevance.

\section*{Summary}

\begin{itemize}
  \item Random assignment identifies the ITT effect of an offer even when
        treatment receipt is endogenous.
  \item With exclusion, relevance, and monotonicity, the outcome ITT divided by
        the first-stage ITT identifies LATE for compliers.
  \item IV is a ratio estimator based on the moment condition
        $\E[Z_i\varepsilon_i]=0$; OLS fails when
        $\Cov(D_i,\varepsilon_i)\neq0$.
  \item Weak first stages make conventional normal approximations unreliable.
  \item A credible instrument requires an institutional argument, not merely a
        significant first stage.
\end{itemize}
